\documentclass[12pt]{article}
\usepackage[spanish]{babel}
\usepackage{amsmath, amssymb, amsthm}
\usepackage{geometry}
\geometry{a4paper, margin=1.5in}
\usepackage{enumitem}
\usepackage{fancyhdr}
\usepackage{hyperref}

% Configuración de encabezado
\pagestyle{fancy}
\fancyhf{}
\rhead{\textbf{Soluciones de Integración}}
\lhead{\thepage}

% Configuración de hipervínculos
\hypersetup{
	colorlinks=true,
	linkcolor=blue,
	filecolor=magenta,      
	urlcolor=cyan,
}

% Comando para problemas
\newcommand{\problema}[1]{\section*{Problema #1}}
\newcommand{\sustituyendo}{\textbf{Sustituyendo:} }
\newcommand{\derivando}{\textbf{Derivando:} }
\newcommand{\integrando}{\textbf{Integrando:} }
\newcommand{\evaluando}{\textbf{Evaluando:} }
\newcommand{\simplificando}{\textbf{Simplificando:} }
\newcommand{\obteniendo}{\textbf{Resultado:} }

\begin{document}
	
	\begin{center}
		\LARGE\textbf{Soluciones de Problemas de Integración} \\
		\large\textbf{Problemas 11-20} \\
		\vspace{0.5cm}
		\normalsize Ejercicios resueltos paso a paso
	\end{center}
	
	\vspace{1cm}
	
	% ========== PROBLEMA 11 ==========
	\problema{11}
	\begin{align*}
		A &= \int_{0}^{\pi/3} \tan x \, dx \\
		\text{Expresando como suma de Riemann:} & \\
		\Delta x &= \frac{\pi/3 - 0}{n} = \frac{\pi}{3n} \\
		x_i &= i\Delta x = \frac{i\pi}{3n} \\
		A &= \lim_{n \to \infty} \sum_{i=1}^{n} \tan\left(\frac{i\pi}{3n}\right) \frac{\pi}{3n}
	\end{align*}
	\textit{Solución mediante sustitución:}
	\begin{align*}
		A &= \int_{0}^{\pi/3} \frac{\sin x}{\cos x} \, dx \\
		\sustituyendo u &= \cos x \\
		\derivando du &= -\sin x \, dx \\
		A &= -\int_{1}^{1/2} \frac{1}{u} \, du \\
		\integrando A &= -\left[\ln|u|\right]_{1}^{1/2} \\
		\evaluando A &= -\left(\ln\frac{1}{2} - \ln 1\right) \\
		\simplificando A &= -\ln\frac{1}{2} = \ln 2 \\
		\obteniendo A &= \boxed{\ln 2}
	\end{align*}
	
	\vspace{0.5cm}
	
	% ========== PROBLEMA 12 ==========
	\problema{12}
	\begin{align*}
		I &= \int_{1}^{3} \left( \frac{1}{2x - 1} - \frac{1}{2x + 1} \right) dx \\
		\sustituyendo 
		I &= \int_{1}^{3} \frac{1}{2x - 1} dx - \int_{1}^{3} \frac{1}{2x + 1} dx
	\end{align*}
	\textit{Resolviendo cada integral por separado:}
	\begin{align*}
		\int_{1}^{3} \frac{1}{2x - 1} dx &= \frac{1}{2} \ln|2x - 1| \Big|_{1}^{3} \\
		\int_{1}^{3} \frac{1}{2x + 1} dx &= \frac{1}{2} \ln|2x + 1| \Big|_{1}^{3} \\
		I &= \frac{1}{2} \left[ \ln|2x - 1| - \ln|2x + 1| \right]_{1}^{3} \\
		\evaluando I &= \frac{1}{2} \left[ (\ln 5 - \ln 7) - (\ln 1 - \ln 3) \right] \\
		\simplificando I &= \frac{1}{2} \left( \ln \frac{5}{7} + \ln 3 \right) \\
		I &= \frac{1}{2} \ln \frac{15}{7} \\
		\obteniendo I &= \boxed{\frac{1}{2} \ln \frac{15}{7}}
	\end{align*}
	
	\vspace{0.5cm}
	
	% ========== PROBLEMA 13 ==========
	\problema{13}
	\begin{align*}
		I &= \int_{2}^{4} \frac{e^{\ln x}}{x^{2} + 7} \, dx \\
		\simplificando I &= \int_{2}^{4} \frac{x}{x^{2} + 7} \, dx \\
		\sustituyendo u &= x^{2} + 7 \\
		\derivando du &= 2x \, dx \\
		I &= \frac{1}{2} \int_{11}^{23} \frac{1}{u} \, du \\
		\integrando I &= \frac{1}{2} \left[ \ln|u| \right]_{11}^{23} \\
		\evaluando I &= \frac{1}{2} (\ln 23 - \ln 11) \\
		\simplificando I &= \frac{1}{2} \ln \frac{23}{11} \\
		\obteniendo I &= \boxed{\frac{1}{2} \ln \frac{23}{11}}
	\end{align*}
	
	\vspace{0.5cm}
	
	% ========== PROBLEMA 14 ==========
	\problema{14}
	\begin{align*}
		I &= \int_{0}^{1} \frac{\arctan x}{x^2 + 1} \, dx \\
		\sustituyendo u &= \arctan x \\
		\derivando du &= \frac{1}{x^2 + 1} \, dx \\
		\text{Cambiando límites: } & x=0 \Rightarrow u=0,\; x=1 \Rightarrow u=\frac{\pi}{4} \\
		I &= \int_{0}^{\pi/4} u \, du \\
		\integrando I &= \left[ \frac{u^2}{2} \right]_{0}^{\pi/4} \\
		\evaluando I &= \frac{1}{2} \left( \frac{\pi}{4} \right)^2 \\
		\simplificando I &= \frac{\pi^2}{32} \\
		\obteniendo I &= \boxed{\frac{\pi^2}{32}}
	\end{align*}
	
	\vspace{0.5cm}
	
	% ========== PROBLEMA 15 ==========
	\problema{15}
	\begin{align*}
		I &= \int_{6}^{9} \frac{3 \ln x - 5}{x} \, dx \\
		\sustituyendo I &= 3 \int_{6}^{9} \frac{\ln x}{x} \, dx - 5 \int_{6}^{9} \frac{1}{x} \, dx
	\end{align*}
	\textit{Resolviendo la primera integral con sustitución:}
	\begin{align*}
		\text{Sea } u &= \ln x \Rightarrow du = \frac{1}{x} dx \\
		\int \frac{\ln x}{x} dx &= \int u \, du = \frac{u^2}{2} = \frac{(\ln x)^2}{2} \\
		\text{Entonces:} \\
		I &= 3 \left[ \frac{(\ln x)^2}{2} \right]{6}^{9} - 5 \left[ \ln x \right]{6}^{9} \\
		\evaluando I &= \frac{3}{2} \left[ (\ln 9)^2 - (\ln 6)^2 \right] - 5(\ln 9 - \ln 6) \\
		\simplificando I &= \frac{3}{2} \left[ (\ln 9)^2 - (\ln 6)^2 \right] - 5 \ln \frac{3}{2} \\
		\obteniendo I &= \boxed{\frac{3}{2} \left[ (\ln 9)^2 - (\ln 6)^2 \right] - 5 \ln \frac{3}{2}}
	\end{align*}
	
	\vspace{0.5cm}
	
	% ========== PROBLEMA 16 ==========
	\problema{16}
	\begin{align*}
		A &= \int_{1}^{3} (x^{3} - 0) \, dx \\
		\simplificando A &= \int_{1}^{3} x^{3} \, dx \\
		\integrando A &= \left[ \frac{x^{4}}{4} \right]_{1}^{3} \\
		\evaluando A &= \frac{3^{4}}{4} - \frac{1^{4}}{4} \\
		\simplificando A &= \frac{81 - 1}{4} = \frac{80}{4} \\
		\obteniendo A &= \boxed{20}
	\end{align*}
	
	\vspace{0.5cm}
	
	% ========== PROBLEMA 17 ==========
	\problema{17}
	\begin{align*}
		A &= \int_{1}^{3} \left[ (4x - x^{2}) - 0 \right] \, dx \\
		\simplificando A &= \int_{1}^{3} (4x - x^{2}) \, dx \\
		\integrando A &= \left[ 2x^{2} - \frac{x^{3}}{3} \right]_{1}^{3} \\
		\evaluando A &= \left(2\cdot 3^2 - \frac{3^3}{3}\right) - \left(2\cdot 1^2 - \frac{1^3}{3}\right) \\
		\simplificando A &= (18 - 9) - \left(2 - \frac{1}{3}\right) \\
		A &= 9 - \frac{5}{3} = \frac{27 - 5}{3} \\
		\obteniendo A &= \boxed{\frac{22}{3}}
	\end{align*}
	
	\vspace{0.5cm}
	
	% ========== PROBLEMA 18 ==========
	\problema{18}
	\begin{align*}
		\text{Para hallar intersecciones:} \\
		1 + y^2 &= 10 \\
		y^2 &= 9 \\
		y &= \pm 3
	\end{align*}
	\begin{align*}
		A &= \int_{-3}^{3} \left[ 10 - (1 + y^2) \right] dy \\
		\simplificando A &= \int_{-3}^{3} (9 - y^2) dy \\
		\integrando A &= \left[ 9y - \frac{y^3}{3} \right]_{-3}^{3} \\
		\evaluando A &= \left(27 - 9\right) - \left(-27 + 9\right) \\
		\simplificando A &= 18 - (-18) = 36 \\
		\obteniendo A &= \boxed{36}
	\end{align*}
	
	\vspace{0.5cm}
	
	% ========== PROBLEMA 19 ==========
	\problema{19}
	\begin{align*}
		\text{Puntos de intersección:} \\
		y^2 + 4y &= 0 \\
		y(y + 4) &= 0 \\
		y &= 0 \quad \text{o} \quad y = -4
	\end{align*}
	\begin{align*}
		A &= \int_{-4}^{0} \left[ 0 - (y^2 + 4y) \right] dy \\
		\simplificando A &= \int_{-4}^{0} (-y^2 - 4y) dy \\
		\integrando A &= \left[ -\frac{y^3}{3} - 2y^2 \right]_{-4}^{0} \\
		\evaluando A &= 0 - \left( -\frac{(-4)^3}{3} - 2(-4)^2 \right) \\
		\simplificando A &= -\left( -\frac{64}{3} - 32 \right) = \frac{64}{3} + 32 \\
		A &= \frac{64 + 96}{3} = \frac{160}{3} \quad \text{(Nota: hay un error en el original)} \\
		\text{Corrigiendo:} \\
		A &= -\left[ \frac{y^3}{3} + 2y^2 \right]_{-4}^{0} \\
		A &= 0 - \left( -\frac{64}{3} + 32 \right) = \frac{64}{3} - 32 \\
		A &= \frac{64 - 96}{3} = -\frac{32}{3} \\
		\text{Valor absoluto del área:} \\
		\obteniendo A &= \boxed{\frac{32}{3}}
	\end{align*}
	
	\vspace{0.5cm}
	
	% ========== PROBLEMA 20 ==========
	\problema{20}
	\begin{align*}
		\text{Área de un sector circular:} \\
		A &= \int_{0}^{a} \frac{1}{2} r^2 \, d\theta \\
		\simplificando A &= \frac{r^2}{2} \int_{0}^{a} d\theta \\
		\integrando A &= \frac{r^2}{2} \left[ \theta \right]_{0}^{a} \\
		\evaluando A &= \frac{r^2}{2} (a - 0) \\
		\obteniendo A &= \boxed{\frac{1}{2} r^2 a}
	\end{align*}


\end{document}